%------------------------------------------------------------------------------%
\documentclass[fleqn,utf8,aspectratio=169,12pt]{beamer}

\usepackage{integracao_e_series}

\title{Séries Numéricas com Termos Não-Negativos}

%------------------------------------------------------------------------------%
\begin{document}

\addtitle

%------------------------------------------------------------------------------%
\section{Série com Termos Não-Negativos}

%------------------------------------------------------------------------------%
\begin{frame}{Recapitulando}
  Sequências numéricas
  \[
    \big( a_k \big)_{k=1}^\infty \qquad\qquad\qquad a_1,\; a_2,\; a_3,\; \dots
  \]
  \pause
  Séries numéricas
  \[
    \sum_{k=1}^\infty a_k = a_1 + a_2 + a_3 + \cdots + a_k + \cdots
  \]
  \pause
  Somas parciais
  \[
    S_n = \sum_{k=1}^n a_k = a_1 + a_2 + a_3 + \cdots + a_n
  \]
\end{frame}

%------------------------------------------------------------------------------%
\begin{frame}{Série com Termos Não-Negativos}
  Caso particular de séries onde \emph{nenhum termo é negativo}
  \pause
  \[
   a_n \geq 0 \quad \forall n
  \]
\end{frame}

%------------------------------------------------------------------------------%
\begin{frame}{Convergência}
  Uma série
  \[
    \sum_{n=1}^\infty a_n \qquad \text{com} \qquad a_n\geq 0
  \]
  \pause
  é \emph{convergente}
  \pause
  se, e somente se, sua sequência de somas parciais, \\
  \pause
  $S_n$, é uma sequência \emph{limitada superiormente}
\end{frame}

%------------------------------------------------------------------------------%
\end{document}
%------------------------------------------------------------------------------%
